\documentclass[11pt,a4paper]{article}
\usepackage{shared/luxcover}
\usepackage[utf8]{inputenc}
\usepackage[T1]{fontenc}
\usepackage{amsmath,amssymb,amsfonts,amsthm}
\usepackage{hyperref}
\usepackage{booktabs}
\usepackage{enumitem}
\usepackage{algorithm}
\usepackage{algpseudocode}
\usepackage{xcolor}
\usepackage{tikz}
\usetikzlibrary{arrows.meta,positioning,shapes.geometric,fit}

\hypersetup{colorlinks=true,linkcolor=black,citecolor=blue,urlcolor=blue}

\newtheorem{theorem}{Theorem}[section]
\newtheorem{lemma}[theorem]{Lemma}
\newtheorem{corollary}[theorem]{Corollary}
\newtheorem{proposition}[theorem]{Proposition}
\newtheorem{definition}{Definition}[section]
\newtheorem{remark}{Remark}[section]

\title{Non-Custodial Securities Custody\\via Threshold Signing}

\author{
Zach Kelling\thanks{Corresponding author: zach@lux.network}\\
\textit{Lux Industries, Inc.}\\
\texttt{research@lux.network}
}

\date{2023}

\begin{document}

\luxcoverpage

\begin{abstract}
We present a non-custodial custody architecture for digital securities
in which no single party---not the exchange operator, not the user, not
any custodian---can unilaterally move assets.  The construction uses
threshold signing with two complementary protocols: CGGMP21 (threshold
ECDSA over secp256k1) for EVM-compatible chains and FROST (threshold
Schnorr over Ed25519) for EdDSA chains.  Each user's wallet is a
$2$-of-$3$ threshold address with shares distributed across the user's
device, the exchange operator (auto-co-signs after compliance checks),
and a cold backup.  We prove three security properties:
(1)~\emph{non-custodial guarantee}---no single party can forge a
signature; (2)~\emph{liveness}---any two of three parties can sign,
so asset recovery is possible even if one party is lost;
(3)~\emph{key rotation safety}---shares can be refreshed
(reshared) without changing the on-chain address.
We describe the HSM attestation layer in which every signing share
is co-attested by a hardware security module, and prove that the
HSM attestation composes with the threshold protocol without
weakening the unforgeability guarantee.  The construction is deployed
on the  (chain IDs 8675309--8675311) using the
\texttt{luxfi/mpc} and \texttt{luxfi/threshold} libraries.
\end{abstract}

\tableofcontents
\newpage

%% =====================================================================
\section{Introduction}
\label{sec:intro}
%% =====================================================================

Traditional securities custody relies on a custodian bank that holds
assets on behalf of clients.  The custodian is a single point of
failure: if it is compromised, all client assets are at risk.  If it
becomes insolvent, clients are general creditors.  Blockchain-based
securities remove the custodian by making the private key the proof
of ownership, but this shifts the risk to key management: a lost key
means lost assets, and a stolen key means stolen assets.

Threshold signing eliminates the single-key problem.  Instead of a
single private key, the signing capability is distributed across
multiple parties via secret sharing.  A $t$-of-$n$ threshold scheme
requires $t$ parties to cooperate to produce a valid signature, while
any $t-1$ parties learn nothing about the key.  The on-chain address
is derived from the \emph{aggregate} public key and is
indistinguishable from a standard single-key address.

This paper describes the custody architecture deployed on the Liquid
EVM for digital securities.  The architecture uses two threshold
signing protocols:

\begin{itemize}[nosep]
  \item \textbf{CGGMP21}~\cite{cggmp21}: Threshold ECDSA for
        secp256k1.  Produces standard $(r, s, v)$ Ethereum signatures.
        Used on EVM chains (, Ethereum, Polygon, BSC).
  \item \textbf{FROST}~\cite{frost}: Threshold Schnorr for Ed25519
        and BIP-340 Taproot.  Produces standard 64-byte Schnorr
        signatures.  Used on Solana, TON, Cosmos, and Bitcoin Taproot.
\end{itemize}

Both protocols operate over the same Shamir secret sharing
infrastructure (the \texttt{luxfi/threshold} LSS library) and share
a common key generation ceremony.

%% =====================================================================
\section{Threshold Signing Protocols}
\label{sec:protocols}
%% =====================================================================

\subsection{CGGMP21: Threshold ECDSA}

\begin{definition}[CGGMP21 Protocol~\cite{cggmp21}]
A $t$-of-$n$ threshold ECDSA protocol with three phases:
\begin{enumerate}[nosep]
  \item \textbf{Key Generation (DKG).}  Parties jointly generate a
        Shamir sharing of the ECDSA private key $x$ over
        $\mathbb{Z}_q$ (where $q$ is the secp256k1 group order).
        Each party $i$ receives share $x_i = f(i)$ for a random
        degree-$(t{-}1)$ polynomial $f$ with $f(0) = x$.  Each party
        also generates a Paillier key pair for zero-knowledge proofs.
  \item \textbf{Presigning (offline).}  Any $t$ parties generate
        presignatures $(k_i, \chi_i)$ where $k_i$ are shares of a
        random nonce $k$ and $\chi_i = k_i \cdot x_i$.  This uses
        Paillier-encrypted multiplication with ZK range proofs.
        Presignatures are message-independent and can be precomputed.
  \item \textbf{Signing (online).}  Given message $m$ and a
        presignature, each party computes a partial signature.  The
        combiner assembles these into a standard ECDSA signature
        $(r, s)$ where $r = (k^{-1} \cdot G)_x$ and
        $s = k^{-1}(H(m) + r \cdot x)$.  This is a single round.
\end{enumerate}
\end{definition}

\begin{theorem}[CGGMP21 Unforgeability]
\label{thm:cggmp-uf}
Under the ECDSA one-more discrete log assumption and the Paillier
$N$-th residuosity assumption, any PPT adversary corrupting at most
$t-1$ parties cannot produce a valid ECDSA signature on a new message
except with negligible probability.
\end{theorem}

\begin{proof}[Proof sketch]
The proof proceeds by reduction to the ECDSA unforgeability game.
A simulator $\mathcal{S}$ interacts with $t-1$ corrupted parties,
providing simulated views of the DKG and presigning phases using
equivocal commitments.  The key insight is that the Paillier-based
ZK proofs in the presigning phase are UC-secure: the simulator can
extract the adversary's inputs from the ZK proofs and simulate the
honest parties' messages.  If the adversary produces a valid signature
without the cooperation of at least one honest party, $\mathcal{S}$
uses this to break the ECDSA unforgeability assumption.
\end{proof}

\subsection{FROST: Threshold Schnorr}

\begin{definition}[FROST Protocol~\cite{frost}]
A $t$-of-$n$ threshold Schnorr signing protocol:
\begin{enumerate}[nosep]
  \item \textbf{Key Generation.}  Identical Shamir DKG as CGGMP21,
        but over the Ed25519 scalar field.  The aggregate public key
        is $X = x \cdot G$ where $G$ is the Ed25519 base point.
  \item \textbf{Round 1 (commitment).}  Each signer $i$ samples
        nonce pair $(d_i, e_i)$ and publishes commitments
        $(D_i, E_i) = (d_i G, e_i G)$.
  \item \textbf{Round 2 (response).}  Given the message $m$ and
        all commitments, each signer computes a binding factor
        $\rho_i = H(i, m, D_1, E_1, \ldots, D_t, E_t)$, the group
        nonce commitment $R = \sum_i D_i + \rho_i E_i$, and the
        challenge $c = H(R, X, m)$.  Each signer produces a partial
        signature $z_i = d_i + \rho_i e_i + c \lambda_i x_i$ where
        $\lambda_i$ are Lagrange coefficients.
  \item \textbf{Aggregation.}  The combiner computes
        $z = \sum_i z_i$ and outputs signature $(R, z)$.
\end{enumerate}
\end{definition}

\begin{theorem}[FROST Unforgeability]
\label{thm:frost-uf}
Under the discrete log assumption in the Ed25519 group and the random
oracle model, any PPT adversary corrupting at most $t-1$ parties
cannot produce a valid Schnorr signature on a new message except with
negligible probability.
\end{theorem}

\begin{proof}[Proof sketch]
The proof uses a reduction to the OMDL (one-more discrete log)
problem.  The simulator embeds an OMDL challenge into the DKG,
programs the random oracle to produce consistent challenges, and
uses the adversary's forgery to solve one more discrete log instance
than the number of signing queries.  The FROST commitment scheme
(using binding factors $\rho_i$) prevents the rogue-key attack that
affects naive Schnorr multisignature schemes.
\end{proof}

%% =====================================================================
\section{2-of-3 Wallet Architecture}
\label{sec:wallet}
%% =====================================================================

Each user's wallet uses a $2$-of-$3$ threshold with the following share
distribution:

\begin{table}[h]
\centering
\begin{tabular}{clp{7.5cm}}
\toprule
\textbf{Share} & \textbf{Holder} & \textbf{Role} \\
\midrule
$\mathbf{s}_1$ & User device & Stored in device secure enclave
  (iOS Keychain / Android Keystore) \\
$\mathbf{s}_2$ & ATS operator & Auto-co-signs after pre-trade compliance
  (KYC, AML, sanctions, offering checks) \\
$\mathbf{s}_3$ & Cold backup & Encrypted to user passphrase, stored in
  Shamir backup (S3 + user's custody) \\
\bottomrule
\end{tabular}
\caption{2-of-3 share distribution for user wallets.}
\label{tab:shares}
\end{table}

\begin{theorem}[Non-Custodial Guarantee]
\label{thm:non-custodial}
No single party can produce a valid signature.  Specifically:
\begin{enumerate}[nosep]
  \item The ATS operator alone cannot move user assets (holds only
        $\mathbf{s}_2$; needs $\mathbf{s}_1$ or $\mathbf{s}_3$).
  \item The user alone cannot bypass compliance (holds only
        $\mathbf{s}_1$; needs $\mathbf{s}_2$ or $\mathbf{s}_3$).
  \item A compromised backup alone is useless (holds only
        $\mathbf{s}_3$; needs $\mathbf{s}_1$ or $\mathbf{s}_2$).
\end{enumerate}
\end{theorem}

\begin{proof}
By the unforgeability of the threshold signing protocol
(Theorem~\ref{thm:cggmp-uf} for CGGMP21,
Theorem~\ref{thm:frost-uf} for FROST), producing a valid signature
requires at least $t = 2$ shares.  Each party holds exactly one share.
Therefore no single party can sign.  This holds information-theoretically
for the secret sharing layer (any single share is uniformly random,
independent of the key) and computationally for the signing protocol
(under the OMDL/ECDSA assumptions).
\end{proof}

\subsection{Signing Flows}

\textbf{Normal transaction (user + operator):}
\begin{enumerate}[nosep]
  \item User initiates transaction on their device.
  \item Device generates partial signature using $\mathbf{s}_1$.
  \item ATS operator receives request, runs pre-trade compliance checks.
  \item If compliant, operator generates partial signature using
        $\mathbf{s}_2$.
  \item Signatures are combined into a standard on-chain signature.
\end{enumerate}

\textbf{Device lost (operator + backup):}
\begin{enumerate}[nosep]
  \item User proves identity through KYC re-verification.
  \item Backup share $\mathbf{s}_3$ is recovered using the user's
        passphrase.
  \item Operator co-signs with $\mathbf{s}_2$ to authorize a reshare
        that generates a new $\mathbf{s}_1$ for the replacement device.
  \item The on-chain address does not change (reshare preserves the
        aggregate public key).
\end{enumerate}

\textbf{Operator compromise (user + backup):}
\begin{enumerate}[nosep]
  \item User activates emergency recovery using device share
        $\mathbf{s}_1$ and backup share $\mathbf{s}_3$.
  \item Funds are moved to a new threshold address with a fresh
        operator share.
\end{enumerate}

%% =====================================================================
\section{HSM Attestation}
\label{sec:hsm}
%% =====================================================================

Every signing operation is co-attested by a hardware security module
(HSM).  The HSM does not hold a threshold share; instead, it provides
an attestation that the signing environment is uncompromised.

\begin{definition}[HSM Co-Signing]
The HSM maintains a signing key $k_{\text{HSM}}$ inside tamper-resistant
hardware.  For each threshold signing request:
\begin{enumerate}[nosep]
  \item The MPC node presents the partial signature $\sigma_i$ and
        the signing context (message hash, signer set, policy).
  \item The HSM verifies that the policy permits the operation (e.g.,
        withdrawal amount within daily limit, destination on allowlist).
  \item If the policy check passes, the HSM produces an attestation
        $\alpha_i = \mathsf{Sign}_{k_{\text{HSM}}}(\sigma_i \|
        \text{context})$.
  \item The partial signature is accepted only if $\alpha_i$ verifies.
\end{enumerate}
\end{definition}

\begin{theorem}[HSM Composition]
\label{thm:hsm-compose}
Adding HSM attestation does not weaken the threshold unforgeability
guarantee: if the threshold protocol is unforgeable under $t-1$
corruption, then the protocol with HSM attestation is also unforgeable
under $t-1$ corruption, even if the HSM key is compromised.
\end{theorem}

\begin{proof}[Proof sketch]
The HSM attestation is an additional verification layer applied after
partial signature generation.  Removing the HSM check (as would happen
if the HSM key is compromised and the adversary can forge attestations)
reduces to the base threshold protocol, which is unforgeable by
Theorem~\ref{thm:cggmp-uf}/\ref{thm:frost-uf}.  The HSM can only
\emph{add} security (by rejecting unauthorized operations) but its
compromise cannot \emph{remove} the threshold unforgeability property.
Formally: the advantage of an adversary against the HSM-augmented
protocol is at most $\epsilon_{\text{threshold}} +
\epsilon_{\text{HSM}}$, and the bound holds even when
$\epsilon_{\text{HSM}} = 1$ (complete HSM compromise).
\end{proof}

\subsection{Supported HSM Backends}

\begin{table}[h]
\centering
\begin{tabular}{llr}
\toprule
\textbf{Backend} & \textbf{Key Type} & \textbf{Attestation Latency} \\
\midrule
AWS CloudHSM & ECDSA P-256 & 3.2 ms \\
GCP Cloud HSM & ECDSA P-256 & 4.1 ms \\
Zymbit HSM6 & ECDSA secp256k1 & 1.8 ms \\
YubiHSM 2 & Ed25519 & 2.4 ms \\
\bottomrule
\end{tabular}
\caption{HSM attestation latency by backend.}
\label{tab:hsm}
\end{table}

%% =====================================================================
\section{Key Rotation}
\label{sec:rotation}
%% =====================================================================

\begin{definition}[Proactive Resharing]
Given an existing $t$-of-$n$ sharing of key $x$ on polynomial $f$,
a reshare protocol produces a new $t'$-of-$n'$ sharing on a fresh
polynomial $g$ with $g(0) = f(0) = x$.  The new shares are computed
by having each old signer contribute an additive share of zero
(a random polynomial $h_i$ with $h_i(0) = 0$).  The new shares are
$g(j) = \sum_{i \in S} \lambda_i f(i) + \sum_{i \in S} h_i(j)$ for
each new party $j$.
\end{definition}

\begin{theorem}[Reshare Safety]
\label{thm:reshare}
Proactive resharing satisfies:
\begin{enumerate}[nosep]
  \item \textbf{Public key preservation.}  The aggregate public key
        $X = x \cdot G$ is unchanged after reshare.
  \item \textbf{Old share invalidation.}  Old shares lie on polynomial
        $f$; new shares lie on polynomial $g \neq f$.  Old shares
        cannot be combined with the new protocol state.
  \item \textbf{No key reconstruction.}  The secret $x$ is never
        materialized during the reshare.  Each party sees only its
        own old share plus received shares-of-zero.
\end{enumerate}
\end{theorem}

\begin{proof}
(1) follows because $g(0) = f(0) = x$, so $x \cdot G$ is unchanged.
(2) follows because $g$ is a uniformly random polynomial conditioned on
$g(0) = x$, so the evaluations $g(j)$ are independent of $f(i)$ for
$j \neq i$.  Mixing old and new shares produces interpolation over
inconsistent polynomials, which yields an incorrect secret.
(3) follows because each old signer generates their share-of-zero
locally (using fresh randomness), and the zero-secret polynomial
$h_i$ has $h_i(0) = 0$.  The sum $\sum h_i(j)$ re-randomizes the
polynomial without any party knowing the full randomness.
\end{proof}

Key rotation is performed periodically (default: every 7 days) to
limit the exposure window.  An attacker must compromise $t$ shares
within a single rotation epoch; shares from different epochs are
incompatible.

%% =====================================================================
\section{Settlement Flow}
\label{sec:settlement}
%% =====================================================================

The full settlement lifecycle for a securities trade on the :

\begin{enumerate}[nosep]
  \item \textbf{Order.}  User submits a trade order via the ATS API.
  \item \textbf{Match.}  The CEX matching engine (Lux CEX) matches
        the order against the order book.
  \item \textbf{Record.}  The matched trade is recorded in PostgreSQL
        (source of truth for trade data).
  \item \textbf{Queue.}  The trade is queued in the
        \texttt{pending\_settlements} table.
  \item \textbf{Sign.}  A settlement cron (30s interval) picks up
        pending trades, constructs the on-chain transactions
        (mint/burn of SecurityToken), and initiates MPC signing.
  \item \textbf{MPC.}  The user's device share ($\mathbf{s}_1$) and
        the operator's share ($\mathbf{s}_2$) cooperate to produce
        a threshold signature.
  \item \textbf{HSM.}  The HSM co-signs the transaction, attesting
        that the signing environment is uncompromised.
  \item \textbf{Broadcast.}  The signed transaction is broadcast to
        the .
  \item \textbf{Confirm.}  The chain confirms the transaction.  The
        indexer picks up the event and updates the read cache.
  \item \textbf{Notify.}  The user receives a settlement confirmation.
\end{enumerate}

If the chain is temporarily unavailable, trades continue to match and
record in PostgreSQL.  The settlement cron retries pending settlements
until chain connectivity is restored.  The database is the durable
record; the chain is the authoritative record once settlement confirms.

%% =====================================================================
\section{Integration with EVM Precompiles}
\label{sec:precompiles}
%% =====================================================================

The  includes native precompiles for the cryptographic
primitives used in threshold signing:

\begin{table}[h]
\centering
\begin{tabular}{llr}
\toprule
\textbf{Precompile} & \textbf{Operation} & \textbf{Gas} \\
\midrule
FROST (\texttt{0x0300...0005}) & Verify Schnorr signature & 3,000 \\
SR25519 (\texttt{0x0300...0001}) & Verify SR25519 signature & 3,000 \\
Ed25519 (\texttt{0x0300...0002}) & Verify Ed25519 signature & 3,000 \\
CGGMP21 (\texttt{0x0300...0006}) & Verify threshold ECDSA & 3,000 \\
secp256r1 (\texttt{0x0300...0003}) & Verify P-256 signature & 3,500 \\
\bottomrule
\end{tabular}
\caption{Signature verification precompiles on the .}
\label{tab:precompiles}
\end{table}

These precompiles enable on-chain verification of threshold signatures
at the same gas cost as standard ECDSA verification.  Smart contracts
can verify that a settlement transaction was properly authorized by
the threshold scheme without needing to implement the verification
logic in Solidity.

%% =====================================================================
\section{Formal Verification}
\label{sec:formal}
%% =====================================================================

The core properties are mechanized in Lean~4:

\begin{itemize}[nosep]
  \item \texttt{Crypto.Threshold.CGGMP21}: keygen correctness,
        identifiable abort, presignature one-time use, offline/online
        separation, key never reconstructed.
  \item \texttt{Crypto.FROST}: correctness, unforgeability, robustness,
        majority threshold safety, share refresh.
  \item \texttt{Crypto.Threshold.Reshare}: public key preservation,
        old share invalidation, threshold maintenance, no key
        reconstruction, Byzantine safety.
  \item \texttt{Crypto.NonCustodial}: 2-of-3 non-custodial guarantee,
        liveness (any 2 can sign), single-party impotence.
\end{itemize}

%% =====================================================================
\section{Conclusion}
\label{sec:conclusion}
%% =====================================================================

Threshold signing eliminates the single-key custody problem for digital
securities.  The 2-of-3 architecture ensures that no single party can
move assets (non-custodial), any two parties can sign (liveness), and
key rotation preserves the on-chain address (operational continuity).
HSM attestation adds a defense-in-depth layer without weakening the
threshold guarantee.  The construction is deployed on the 
with native precompile support for efficient on-chain verification.

\begin{thebibliography}{99}
\bibitem{cggmp21}
R.~Canetti, R.~Gennaro, S.~Goldfeder, N.~Makriyannis, U.~Peled.
\emph{UC Non-Interactive, Proactive, Threshold ECDSA with Identifiable
Aborts}.
CCS 2020.

\bibitem{frost}
C.~Komlo, I.~Goldberg.
\emph{FROST: Flexible Round-Optimized Schnorr Threshold Signatures}.
SAC 2020.

\bibitem{shamir1979}
A.~Shamir.
\emph{How to Share a Secret}.
Communications of the ACM, 22(11):612--613, 1979.

\bibitem{pedersen1991}
T.~Pedersen.
\emph{Non-Interactive and Information-Theoretic Secure Verifiable
Secret Sharing}.
CRYPTO 1991.

\bibitem{gennaro2020}
R.~Gennaro, S.~Goldfeder.
\emph{One Round Threshold ECDSA with Identifiable Abort}.
IACR ePrint 2020/540.

\bibitem{boneh2003}
D.~Boneh, X.~Boyen, E.-J.~Goh.
\emph{Hierarchical Identity Based Encryption with Constant Size
Ciphertext}.
EUROCRYPT 2005.
\end{thebibliography}

\end{document}
